% 2026-11-11 / v39
\section{混合cubicのpoint schemeの非被約構造}
\label{ps39:sec}
以下では混合族$\lambda=1$、$h\in\mathbb C^\times$、
$h^6\ne-27$を固定する。$B_h=A_h/(\Omega_h)$とし、
point schemeとはHilbert級数$(1-t)^{-1}$のgraded右point moduleの
初期点を表す射影スキームを意味する。
スキーム論的な計算と台集合の計算を区別する。

\subsection{中心一次元の非零チャート}
\begin{lemma}[中心座標による点列の対角化]\label{ps39:lem:abelian}
$u\ne0$におけるpoint moduleの族は、スキームとして
$B_h$の可換化の$u\ne0$チャートと一致する。
\end{lemma}
\begin{proof}
可換試験環$T$上の点列を$p_i=(x_i:y_i:z_i:u_i:v_i)$と表す。
多重線形化した中心性$[u,a]=0$（任意の一次生成元$a$）から
$u_ia_{i+1}-a_iu_{i+1}=0$を得る。$u_i$が単元なら各点を
$u_i=1$に規格化でき、$a_{i+1}=a_i$となる。
全$i$について繰り返すと点列はスキーム論的に一定である。
残る多重線形化関係式は、すべて同一座標への可換代入に等しい。
逆に可換化の任意の$T$値点は一定点列を定義する。
この対応は$T$に自然なので主張が従う。
\end{proof}

\begin{proposition}[孤立部分は12個の二重点]\label{ps39:prop:length}
$u=1$におけるpoint schemeのイデアルは
\begin{equation}\label{ps39:eq:ideal}
I_h=\left(x^2+y,\ y^2+x,\ z^2,\
1+v^3+\frac{xy}{3}
+h^2z\left(\frac16-\frac{2xy}{3}\right)\right).
\end{equation}
これは12個の相異なる点に支持を持つ長さ2の非被約点の非交和で、
総長は24である。
\end{proposition}
\begin{proof}
$h\ne0$なので交換子を可換化すると
$x^2+uy=y^2+ux=z^2=0$となる。
第\ref{mt:sec:family}節の中心三次元の可換化は
\[
\Omega_h^{\mathrm{ab}}
=F_1-\frac{2h^2}{3}xyz+\frac{h^2u^2}{6}z
+\frac{h^4u}{18}z^2.
\]
$u=1$で最初の関係を使うと$x^3=y^3=-xy$、$z^3=0$であり、
この式が\eqref{ps39:eq:ideal}へ変わる。
$x^2+y=y^2+x=0$は$(x,y)=(0,0)$、または
$x^3=-1,\ y=-x^2$という計四個の解を持つ。
Jacobian行列式は$4xy-1$で、前者で$-1$、後者で$3$だから
すべて被約である。$xy=0$または$1$に従い、
$z=0$で$v^3=-1$または$-4/3$となる。
いずれも三つの相異なる非零根を持ち、支持は計12点。
各点で$3v^2$が単元なので最後の式は$v$を$z$の冪零変形として
一意に解く。残る条件は$z^2=0$のみであり、
各局所環は$\mathbb C[z]/(z^2)$に同型である。
\end{proof}

\subsection{Poisson零点との同型と埋め込みの相違}
\begin{theorem}[孤立部分の比較]\label{ps39:thm:comparison}
$X_0=V(F_1)$上のPoisson零点の孤立部分を$Z_{\mathrm{crit}}$、
非可換point schemeの孤立部分を$\Sigma_h$とする。
このとき$\Sigma_h\simeq Z_{\mathrm{crit}}$がスキームとして成立するが、
$h\ne0$では両者は$\mathbb P^4$内の同じ閉部分スキームではない。
$u=1$での同型は
\begin{equation}\label{ps39:eq:shear}
(x,y,z,v)\longmapsto
\left(x,y,z,\
v'=v+\frac{h^2z}{3v^2}
 \left(\frac16-\frac{2xy}{3}\right)\right)
\end{equation}
で与えられる。
\end{theorem}
\begin{proof}
Poisson零点の孤立部分は
\[
I_0=(x^2+y,\ y^2+x,\ z^2,\ 1+v^3+xy/3)
\]
で定義される。両スキーム上で$v^3=-1$または$-4/3$を法とする
冪零元の補正だけがあるので、$v$は単元である。
$z^2=0$より
\[
(v')^3=v^3+h^2z(1/6-2xy/3).
\]
従って\eqref{ps39:eq:shear}で$I_h$と$I_0$が対応し、
逆写像は$z$の一次のせん断を逆向きに施すことで得られる。

一方、$(x,y,z)=(0,0,0)$、$v_0^3=-1$の点で
Poisson零点の埋め込み接方向は$\delta v=0$、
point schemeの接方向は
$3v_0^2\delta v+(h^2/6)\delta z=0$である。
$h\ne0$では異なるため、閉部分スキームとしては一致しない。
\end{proof}

\subsection{Sklyanin曲線と全point scheme}
\begin{theorem}[一般非零パラメータのpoint scheme]\label{ps39:thm:global}
$h\ne0$、$h^6\ne-27$のとき、point moduleの初期点を表す
スキームは
\[
\boxed{\quad
\operatorname{Point}(B_h)
\simeq E_h\sqcup\coprod_{a=1}^{12}
\operatorname{Spec}\mathbb C[\varepsilon]/(\varepsilon^2).
\quad}
\]
ここで$E_h=V(u,v,x^3+y^3+z^3+h^2xyz)$は滑らかな楕円曲線である。
特に$h=0$のpoint schemeは可換cubic $X_0$全体であって、
上の表示が$h=0$までpoint scheme全体の平坦族をなすとはしない。
\end{theorem}
\begin{proof}
$u\ne0$では直前の命題を使う。
$u=0,v\ne0$なら中心性から点列は一定であるが、
可換化した交換子は$x=y=z=0$を強制し、
中心三次元の値は$v^3\ne0$なのでpoint moduleは存在しない。
残りの支持は$u=v=0$であり、三生成元Sklyanin代数
\[
S_h=\mathbb C\langle x,y,z\rangle/
(yz-zy-hx^2,\ zx-xz-hy^2,\ xy-yx-hz^2)
\]
のpoint schemeとなる。既知のSklyanin定理
\cite{PSPSmithTate,PSPDeLaet}によりこれは滑らかな$E_h$で、
その中心三次元はすべてのpoint moduleを消す。
$\Omega_h\bmod(u,v)$はその中心三次元の非零倍なので、
これらのpoint moduleは$B_h$へ降りる。

さらに$E_h$方向に$u,v$の冪零的な厚化が生じないことを示す。
Sklyaninのshiftは$E_h$上の平行移動で、$h\ne0$なら恒等ではない。
従って$p_0$と$p_1$は$E_h$のどの点でも異なる。
その点の局所環である二次小行列
$p_{0,a}p_{1,b}-p_{0,b}p_{1,a}$が単元になる。
中心性の多重線形化
$u_0p_{1,a}=p_{0,a}u_1$および
$v_0p_{1,a}=p_{0,a}v_1$を二座標について消去すると、
その小行列と$u_0,v_0$との積は零となる。
従って局所環上で$u_0=v_0=0$が厳密に成立し、
近傍のpoint schemeはSklyaninの被約$E_h$に等しい。
$E_h$と$u\ne0$部分は互いに交わらないので非交和となる。
\end{proof}

\begin{remark}[結論の範囲]
Poisson零点全体$E_0\sqcup Z_{\mathrm{crit}}$と
非可換point scheme$E_h\sqcup\Sigma_h$は一般には一致しない。
孤立部分は抽象スキームとして同型だが、法線方向の埋め込みは異なる。
復元cubicの$u=v=0$制限は$E_h$に等しいものの、
point scheme全体が復元cubic上にあるとは主張しない。
Sklyanin退化値$h^6=-27$と$h=0$は上の定理の範囲外である。
\end{remark}
